% GENERATED FILE. Edit canonical lessons, then run npm run generate:books.
% canonical-source-hash: fnv1a64:7b698ba79a5dc72f
% canonical-lessons: IT-C340-dibattito, IT-C340-tesi, IT-C340-argomentazione, IT-C340-convincente, IT-C340-persuadere

\chapter{Il dibattito}
\label{ch:it-il-dibattito}
\hlchaptermodality{340}

\begin{chapteropening}
\textbf{By the end of this chapter:} \emph{I can take part in a debate: state a thesis, build an argument and try to persuade with a convincing reason.}

Complete the last lesson of chapter 340: I can take part in a debate: state a thesis, build an argument and try to persuade with a convincing reason.
\end{chapteropening}

\section[il dibattito]{il dibattito — the debate}

\label{lesson:IT-C340-dibattito}



\begin{quote}
Before the new one: say the Italian for to maintain, then the Italian for the point of view.
\end{quote}

\subsection*{You'll want to know: il dibattito}
\textbf{il dibattito} — \textquotedblleft{}the debate\textquotedblright{}.

\begin{quote}
Stasera c'è un dibattito in televisione sulla scuola.
\end{quote}

\begin{quote}
\emph{Tonight there is a debate on television about schools.}
\end{quote}

\emph{Un dibattito acceso} — a heated debate.

\begin{cousinweb}
\emph{dibattito} is built on \emph{dibattere} \textquotedblleft{}to debate\textquotedblright{}, from \emph{dis-} and \emph{battere} \textquotedblleft{}to beat\textquotedblright{}, which is inherited from Latin \emph{battuere} \textquotedblleft{}to beat, to strike\textquotedblright{} (Lewis \& Short). Ideas are knocked back and forth.
\end{cousinweb}

\begin{grammarlens}[title={What you've built}]
One more word for telling, explaining and describing at B1.
\end{grammarlens}

\subsection*{Your turn}
\begin{itemize}
\raggedright
  \item \emph{Say it:} \emph{il dibattito}
  \item \emph{Say it:} \emph{il dibattito}, once more
  \item \emph{Say it:} a sentence of your own with \emph{il dibattito}
  \item \emph{Recall:} say the Italian for to maintain, then the Italian for the point of view, then say \emph{il dibattito} again
\end{itemize}

\begin{tcolorbox}[breakable,colback=teal!4,colframe=teal!35!black,title={Before you move on}]
What does il dibattito mean? (\textquotedblleft{}The debate\textquotedblright{}.) Say it once more.
\end{tcolorbox}
\section[la tesi]{la tesi — the thesis, the position (argued for)}

\label{lesson:IT-C340-tesi}



\begin{quote}
Before the new one: say the Italian for the point of view, then the Italian for the debate.
\end{quote}

\subsection*{You'll want to know: la tesi}
\textbf{la tesi} — \textquotedblleft{}the thesis, the position (argued for)\textquotedblright{}.

\begin{quote}
La sua tesi è semplice: lavoriamo troppo.
\end{quote}

\begin{quote}
\emph{His thesis is simple: we work too much.}
\end{quote}

The same in the plural: \emph{le tesi}. At university, \emph{la tesi} is the dissertation.

\begin{cousinweb}
\emph{tesi} is a learned borrowing from Latin \emph{thesis} \textquotedblleft{}a proposition\textquotedblright{} (Lewis \& Short), which came from Greek \emph{thésis} \textquotedblleft{}a setting down\textquotedblright{}, from the verb \emph{tithénai} \textquotedblleft{}to put\textquotedblright{}.
\end{cousinweb}

\begin{grammarlens}[title={What you've built}]
One more word for telling, explaining and describing at B1.
\end{grammarlens}

\subsection*{Your turn}
\begin{itemize}
\raggedright
  \item \emph{Say it:} \emph{la tesi}
  \item \emph{Say it:} \emph{la tesi}, once more
  \item \emph{Say it:} a sentence of your own with \emph{la tesi}
  \item \emph{Recall:} say the Italian for the point of view, then the Italian for the debate, then say \emph{la tesi} again
\end{itemize}

\begin{tcolorbox}[breakable,colback=teal!4,colframe=teal!35!black,title={Before you move on}]
What does la tesi mean? (\textquotedblleft{}The thesis, the position (argued for)\textquotedblright{}.) Say it once more.
\end{tcolorbox}
\section[l'argomentazione]{l'argomentazione — the line of argument, the reasoning}

\label{lesson:IT-C340-argomentazione}



\begin{quote}
Before the new one: say the Italian for the debate, then the Italian for the thesis.
\end{quote}

\subsection*{You'll want to know: l'argomentazione}
\textbf{l'argomentazione} — \textquotedblleft{}the line of argument, the reasoning\textquotedblright{}.

\begin{quote}
La tua argomentazione è chiara, ma non sono d'accordo.
\end{quote}

\begin{quote}
\emph{Your line of argument is clear, but I do not agree.}
\end{quote}

The chain of reasons you give; one single reason is \emph{un argomento}.

\begin{cousinweb}
\emph{argomentazione} is a learned borrowing from Latin \emph{argūmentātiō} \textquotedblleft{}the bringing of proof\textquotedblright{} (Lewis \& Short), from \emph{argūmentum} \textquotedblleft{}proof, evidence\textquotedblright{}, built on \emph{arguere} \textquotedblleft{}to make clear, to prove\textquotedblright{}.
\end{cousinweb}

\begin{grammarlens}[title={What you've built}]
One more word for telling, explaining and describing at B1.
\end{grammarlens}

\subsection*{Your turn}
\begin{itemize}
\raggedright
  \item \emph{Say it:} \emph{l'argomentazione}
  \item \emph{Say it:} \emph{l'argomentazione}, once more
  \item \emph{Say it:} a sentence of your own with \emph{l'argomentazione}
  \item \emph{Recall:} say the Italian for the debate, then the Italian for the thesis, then say \emph{l'argomentazione} again
\end{itemize}

\begin{tcolorbox}[breakable,colback=teal!4,colframe=teal!35!black,title={Before you move on}]
What does l'argomentazione mean? (\textquotedblleft{}The line of argument, the reasoning\textquotedblright{}.) Say it once more.
\end{tcolorbox}
\section[convincente]{convincente — convincing}

\label{lesson:IT-C340-convincente}



\begin{quote}
Before the new one: say the Italian for the thesis, then the Italian for the line of argument.
\end{quote}

\subsection*{You'll want to know: convincente}
\textbf{convincente} — \textquotedblleft{}convincing\textquotedblright{}.

\begin{quote}
Non è una risposta molto convincente.
\end{quote}

\begin{quote}
\emph{It is not a very convincing answer.}
\end{quote}

Opposite: \emph{poco convincente}. One form for both genders.

\begin{cousinweb}
\emph{convincente} is the present participle of \emph{convincere}, which comes from Latin \emph{convincere} \textquotedblleft{}to overcome in argument, to prove\textquotedblright{} (Lewis \& Short), built on \emph{vincere} \textquotedblleft{}to conquer\textquotedblright{}.
\end{cousinweb}

\begin{grammarlens}[title={What you've built}]
One more word for telling, explaining and describing at B1.
\end{grammarlens}

\subsection*{Your turn}
\begin{itemize}
\raggedright
  \item \emph{Say it:} \emph{convincente}
  \item \emph{Say it:} \emph{convincente}, once more
  \item \emph{Say it:} a sentence of your own with \emph{convincente}
  \item \emph{Recall:} say the Italian for the thesis, then the Italian for the line of argument, then say \emph{convincente} again
\end{itemize}

\begin{tcolorbox}[breakable,colback=teal!4,colframe=teal!35!black,title={Before you move on}]
What does convincente mean? (\textquotedblleft{}Convincing\textquotedblright{}.) Say it once more.
\end{tcolorbox}
\section[persuadere]{persuadere — to persuade}

\label{lesson:IT-C340-persuadere}



\begin{quote}
Before the new one: say the Italian for the line of argument, then the Italian for convincing.
\end{quote}

\subsection*{You'll want to know: persuadere}
\textbf{persuadere} — \textquotedblleft{}to persuade\textquotedblright{}.

\begin{quote}
Ho cercato di persuadere mio padre a comprare una macchina nuova.
\end{quote}

\begin{quote}
\emph{I tried to persuade my father to buy a new car.}
\end{quote}

\emph{Persuadere qualcuno a fare qualcosa}. Past participle: \emph{persuaso}.

\begin{cousinweb}
\emph{persuadere} is a learned borrowing from Latin \emph{persuādēre} \textquotedblleft{}to bring over by talking, to convince\textquotedblright{} (Lewis \& Short), built from \emph{per-} and \emph{suādēre} \textquotedblleft{}to advise, to urge\textquotedblright{}.
\end{cousinweb}

\begin{grammarlens}[title={What you've built}]
One more word for telling, explaining and describing at B1.
\end{grammarlens}

\subsection*{Your turn}
\begin{itemize}
\raggedright
  \item \emph{Say it:} \emph{persuadere}
  \item \emph{Say it:} \emph{persuadere}, once more
  \item \emph{Say it:} a sentence of your own with \emph{persuadere}
  \item \emph{Recall:} say the Italian for the line of argument, then the Italian for convincing, then say \emph{persuadere} again
\end{itemize}

\begin{tcolorbox}[breakable,colback=teal!4,colframe=teal!35!black,title={Before you move on}]
What does persuadere mean? (\textquotedblleft{}To persuade\textquotedblright{}.) Say it once more.
\end{tcolorbox}
